% !TeX root = ../navier-stokes-manuscript.tex
% ============================================================
% 2.6 Finite Adjacent Rectangles
% ============================================================

\subsection{Finite Adjacent Rectangles}\label{sub:FiniteAdjacentRectangles}

Begin with the first finite case that carries two temporal rows to a common trace height.
Fix a strict interval \(I=(0,T)\), \(T<T_*\), and write
\[
  V_0=u,
  \qquad
  V_1=\partial_tu,
  \qquad
  V_2=\partial_t^2u.
\]
Suppose
\[
  V_0,V_1
  \in
  L^2(I;H^3(\R^3))
\]
and
\[
  V_1,V_2
  \in
  L^2(I;H^1(\R^3)).
\]
For \(V_0\), these assumptions give
\[
  V_0\in L^2(I;H^3),
  \qquad
  \partial_tV_0=V_1\in L^2(I;H^1).
\]
The Hilbert triple
\[
  H^3(\R^3)
  \hookrightarrow
  H^2(\R^3)
  \hookrightarrow
  H^1(\R^3)
\]
uses \(H^2\) as the pivot space.
For a smooth function \(v\), the identity behind its trace is
\[
  \norm{v(t)}_{H^2}^2
  -
  \norm{v(s)}_{H^2}^2
  =
  2\int_s^t
  \left\langle
    \partial_\tau v(\tau),
    v(\tau)
  \right\rangle_{H^1,H^3;H^2}
  \,d\tau.
\]
The right side is integrable when
\[
  v\in L^2(I;H^3),
  \qquad
  \partial_tv\in L^2(I;H^1).
\]
The standard density argument then gives \(v\in C([0,T];H^2)\).
Applying this identity to \(V_0\) gives
\[
  V_0\in C\bigl([0,T];H^2(\R^3)\bigr).
\]
Applying the same argument to \(V_1\), using \(\partial_tV_1=V_2\), gives
\[
  V_1\in C\bigl([0,T];H^2(\R^3)\bigr).
\]
The two adjacent spatial rows have produced time traces for both retained temporal coordinates.

Collect the four norms in
\begin{align*}
  \mathfrak R_{1,2}(u;T)
  :={}&
  \norm{u}_{L^2(0,T;H^3)}^2
  +
  \norm{\partial_tu}_{L^2(0,T;H^3)}^2
  \\
  &+
  \norm{\partial_tu}_{L^2(0,T;H^1)}^2
  +
  \norm{\partial_t^2u}_{L^2(0,T;H^1)}^2.
\end{align*}
Its coordinate pattern is
\[
  \begin{array}{c|cc}
    & n=0 & n=1
    \\
    \hline
    H^3
    & u & \partial_tu
    \\
    H^1
    & \partial_tu & \partial_t^2u
  \end{array}
\]
The top and bottom entries in each column form one adjacent Hilbert-triple pair.
The repeated \(\partial_tu\) joins the first temporal block to the second.
This finite two-by-two pattern is the first adjacent rectangle.

The same construction works at arbitrary finite depths.
Fix independent integers \(N,M\in\Nzero\).
The middle trace height is \(H^M\), with adjacent spaces
\[
  H^{M+1}(\R^3)
  \hookrightarrow
  H^M(\R^3)
  \hookrightarrow
  H^{M-1}(\R^3).
\]
For \(M=0\), the lower space is the inhomogeneous space \(H^{-1}(\R^3)\).
Define the finite velocity norm
\begin{align*}
  \mathfrak R_{N,M}(u;T)
  :={}&
  \sum_{n=0}^{N}
  \norm{V_n}_{L^2(0,T;H^{M+1})}^{2}
  \\
  &+
  \sum_{n=0}^{N}
  \norm{V_{n+1}}_{L^2(0,T;H^{M-1})}^{2}.
\end{align*}
Finiteness of this quantity gives, for each \(0\leq n\leq N\),
\[
  V_n\in L^2(0,T;H^{M+1}),
  \qquad
  \partial_tV_n=V_{n+1}\in L^2(0,T;H^{M-1}).
\]
For a smooth representative, these two rows meet in the exact calculation
\[
  \left|
    \frac{d}{dt}
    \norm{V_n(t)}_{H^M}^{2}
  \right|
  \leq
  2
  \norm{V_{n+1}(t)}_{H^{M-1}}
  \norm{V_n(t)}_{H^{M+1}}.
\]
The right side belongs to \(L^1(0,T)\) by Cauchy--Schwarz.
Approximation by smooth functions, together with weak continuity, therefore gives the \(H^M\)-continuous representative.
The same trace argument yields
\[
  V_n
  \in
  C\bigl([0,T];H^M(\R^3)\bigr),
  \qquad
  0\leq n\leq N.
\]
Increasing \(N\) adds temporal rows.
Increasing \(M\) raises the spatial trace height.
The two choices remain independent.

The number \(\mathfrak R_{N,M}(u;T)\) measures velocity rows only.
The pressure-complete finite record must identify those rows as derivatives of the same Navier--Stokes history.
For the fixed whole-space solution \((u,p)\) generated by \(u_0\), define
\[
  \cR_{N,M}[u_0;T]
  :=
  \left(
    (V_n)_{n=0}^{N+1},
    (Q_n)_{n=0}^{N};
    \mathfrak R_{N,M}(u;T)
  \right),
\]
where
\[
  V_n=\partial_t^nu,
  \qquad
  Q_n=\partial_t^np,
\]
the velocity rows have the memberships counted by \(\mathfrak R_{N,M}\), and for every \(0\leq n\leq N\),
\[
  V_{n+1}
  +
  \sum_{a=0}^{n}
  \binom{n}{a}
  (V_a\cdot\nabla)V_{n-a}
  +\nabla Q_n
  =
  \nu\Delta V_n,
\]
\[
  \nabla\cdot V_n=0,
\]
and
\[
  -\Delta Q_n
  =
  \partial_i\partial_j
  \sum_{a=0}^{n}
  \binom{n}{a}
  (V_a)_i(V_{n-a})_j.
\]
The rapid-decay normalization fixes each \(Q_n\).
Every mixed coordinate \(J_{n,\alpha}\) available at the retained Sobolev depths is the weak spatial derivative \(\partial_x^\alpha V_n\), and its pressure coordinate is the corresponding derivative of \(Q_n\).
Thus \(\mathfrak R_{N,M}\) is the finite velocity norm, while \(\cR_{N,M}\) is the finite pressure-complete Navier--Stokes object built around that norm.

Every finite rectangle supplies only the traces at its chosen temporal and spatial depths.
The all-orders velocity intersection is recovered by allowing all finite pairs \((N,M)\).
Those rectangles belong to one intersection when every shared coordinate agrees, which is the compatibility question.
